\documentclass[11pt,a4paper]{article}
\usepackage[T1]{fontenc}
\usepackage[utf8]{inputenc}
\usepackage{newtxtext,newtxmath}
\usepackage[scaled=0.94]{sourcesanspro}
\usepackage{geometry}
\geometry{left=23mm,right=23mm,top=22mm,bottom=23mm,headheight=16pt}
\usepackage{microtype,parskip,booktabs,longtable,array,tabularx}
\usepackage{amsmath,mathtools}
\usepackage[dvipsnames,table]{xcolor}
\usepackage{enumitem,fancyhdr,titlesec,tcolorbox,listings,needspace}
\usepackage{hyperref,xurl}
\definecolor{DWGreen}{HTML}{214E46}
\definecolor{DWInk}{HTML}{24323B}
\definecolor{DWGold}{HTML}{B18C49}
\definecolor{DWLight}{HTML}{F1F5F3}
\definecolor{DWGray}{HTML}{59666C}
\definecolor{DWLine}{HTML}{D9E2DE}
\hypersetup{colorlinks=true,linkcolor=DWGreen,urlcolor=DWGreen,
 pdftitle={Dawnwood Interactive: Unified Substrate Definition},
 pdfauthor={Tom Klootwijk | Dawnwood Interactive},
 pdfsubject={Universal Spatial State Automaton | NL200678942 | Version 0.2},
 pdfkeywords={Dawnwood Interactive, Tom Klootwijk, NL200678942, log-polar, Klein bottle, SDF, self-referential}}
\color{DWInk}
\setlength{\parskip}{6pt}
\setlength{\parindent}{0pt}
\clubpenalty=10000 \widowpenalty=10000
\setlength{\emergencystretch}{2em}
\renewcommand{\arraystretch}{1.2}
\setlength{\tabcolsep}{5pt}
\setlist{nosep,leftmargin=1.4em,topsep=4pt}
\titleformat{\section}{\Large\sffamily\bfseries\color{DWGreen}}{\thesection}{0.65em}{}
\titleformat{\subsection}{\large\sffamily\bfseries\color{DWInk}}{\thesubsection}{0.65em}{}
\titlespacing*{\section}{0pt}{0pt}{11pt}
\titlespacing*{\subsection}{0pt}{12pt}{5pt}
\pagestyle{fancy}\fancyhf{}
\fancyhead[L]{\sffamily\small\color{DWGreen}DAWNWOOD INTERACTIVE}
\fancyhead[R]{\sffamily\small\color{DWGray}UNIFIED DEFINITION}
\fancyfoot[L]{\sffamily\footnotesize\color{DWGray}Tom Klootwijk \enspace | \enspace Version 0.2}
\fancyfoot[R]{\sffamily\small\thepage}
\renewcommand{\headrulewidth}{0.4pt}
\renewcommand{\footrulewidth}{0pt}
\newcolumntype{Y}{>{\raggedright\arraybackslash}X}
\newcolumntype{P}[1]{>{\raggedright\arraybackslash}p{#1}}
\newcommand{\src}[1]{{\sffamily\footnotesize\color{DWGray}[Source, pp.~#1]}}
\newcommand{\code}[1]{\texttt{#1}}
\newcommand{\Klein}{\mathcal{K}}
\newcommand{\LUT}{\mathcal{L}}
\newcommand{\State}{\mathcal{S}}
\newcommand{\Op}{\mathcal{O}}
\newcommand{\Apply}{\operatorname{Apply}}
\newcommand{\Pin}{\operatorname{Pin}}
\newcommand{\Inv}{\operatorname{Inv}}
\newtcolorbox{dwbox}[1][]{colback=DWLight,colframe=DWLine,boxrule=0.5pt,arc=1mm,
 left=3mm,right=3mm,top=2.5mm,bottom=2.5mm,fonttitle=\sffamily\bfseries,coltitle=DWGreen,title={#1}}
\lstset{basicstyle=\ttfamily\small,breaklines=true,frame=single,framerule=0.3pt,
 rulecolor=\color{DWLine},backgroundcolor=\color{DWLight},columns=fullflexible,
 keepspaces=true,showstringspaces=false,xleftmargin=2pt,xrightmargin=2pt}
\setcounter{tocdepth}{1}
\begin{document}
\begin{titlepage}
\vspace*{8mm}
{\sffamily\bfseries\color{DWGreen}\Large DAWNWOOD INTERACTIVE}\par
\vspace{6mm}{\color{DWGold}\rule{48mm}{1.4pt}}\par
\vspace{17mm}
{\sffamily\bfseries\fontsize{32}{37}\selectfont Unified\\Substrate Definition}\par
\vspace{9mm}
{\sffamily\fontsize{17}{22}\selectfont Self-Referential Log-Polar\\Klein-Bottle SDF}\par
\vspace{5mm}
{\sffamily\color{DWGreen}\large Universal Spatial State Automaton}\par
\vspace{18mm}
\begin{tabular}{@{}ll@{}}
\textsf{Author} & \textbf{Tom Klootwijk}\\[7pt]
\textsf{Identifier} & \textbf{NL200678942}\\[7pt]
\textsf{Date of birth} & \textbf{10-07-1990}\\[7pt]
\textsf{Telephone} & \textbf{+31655954068}\\[7pt]
\textsf{Project} & Dawnwood Interactive\\[7pt]
\textsf{Edition} & Version 0.2 \enspace / \enspace 16 September 2026
\end{tabular}
\vfill
\begin{dwbox}
One circulating definition: the double pinion carries the wavefront and its operators on the Klein-bottle SDF; the log-polar LUT and implicit BST hold their relations; one-bit jitter at interval $\Psi$ changes the living operator field; the returned state becomes its next input.
\end{dwbox}
\vspace{5mm}
{\sffamily\small\color{DWGray}Source: \code{double-slit-theory.pdf}, 24 pages.\par
Editable manuscript, substrate model and symbolic workbench accompany this edition.}
\end{titlepage}
\setcounter{page}{2}
\tableofcontents
\vfill
\begin{dwbox}[Reading the definition]
The document consolidates the source discussion in its own system language. Equations, record names and the written cycle are editable notation for those relationships. A numerical law absent from the source remains a named operator, with its binding exposed in the accompanying model.

The source page references lead back to the unchanged discussion. The author details on the title page follow the details supplied directly for this edition.
\end{dwbox}
\clearpage

\section{The unified Dawnwood definition}
\textbf{Dawnwood Interactive is the definition of a self-referential spatial computing substrate in which the wavefront, executable operators, operator positions, branching relations and recurrent state occupy one Klein-bottle SDF system.} The two vertical slits become the double pinion and Hadamard mitosis hinges. Their dichromatic streams are carried through log-encoded polar coordinates, a geometric LUT and a one-bit BST. At interval $\Psi$, one-bit jitter acts on the operators themselves. The field that acts is also the field that changes. \src{3--16}

The source names this whole a \emph{Universal Spatial State Automaton}. Its expression language places branching, state, transformation, feedback and operator change within the same spatial definition. The source also describes the circulating operators as dynamic phase shifters, stochastic logic gates and self-assembling cellular nodes. These names refer to roles inside the Dawnwood definition. \src{15--16, 22--23}

\subsection{One state, including its operators}
Write the circulating state as
\begin{equation}
\State_n = \bigl(\Klein_n,\, W_n,\,\LUT_n,\,\mathcal{B}_n,\,C_n\bigr),
\end{equation}
where $\Klein$ is the Klein-bottle SDF relation, $W$ the wavefront, $\LUT$ the SDF operator LUT, $\mathcal{B}$ the implicit binary routing structure and $C=(R,G,B,A)$ the texture state. These are coordinates of one definition, rather than separate engines placed beside one another.

Its central recurrence is
\begin{equation}
\boxed{\State_{n+1}=\mathfrak{D}_{\Psi}\bigl[\State_n;j_n\bigr],
\qquad j_n\in\{0,1\}.}
\end{equation}
The operator field $\LUT_n$ is part of the input to $\mathfrak{D}_{\Psi}$ and part of its output. A change to an operator body or surface position therefore changes a subsequent use of the same definition. The output is returned as input through the non-orientable surface relation. \src{13--16}

\subsection{The circulation is the primary object}
The source places rasterization, raymarching, raytracing and square output fields outside this core description. Optional Bayer output is downstream of the circulation. It reads an already-formed state; it does not define the surface on which computation takes place. \src{12--14}

The development object is consequently the whole recurrence: double pinion, phase, jitter, geometry, branching, mutation and return. The supporting workbench keeps their expression records connected throughout each $\Psi$ cycle.
\clearpage

\section{Source progression and common notation}
The source develops its vocabulary as a single continuing construction. The following sequence retains the technical milestones and their clock labels, with the journey conversation left in the original attachment.

{\small
\begin{tabularx}{\linewidth}{@{}P{15mm}YP{16mm}@{}}
\toprule
\textbf{Label} & \textbf{Definition entering the whole} & \textbf{Pages}\\
\midrule
07:52 & Double vertical slit, two pinions, Hadamard mitosis hinges, lowercase phi, log-polar LUT. & 3--4\\
07:54 & 0-order notation, bit shifting and one-bit even/odd parity. & 4--5\\
07:56 & The 756 wavefront FT vector $[0,2,0,1]$ and one-bit jitter. & 6\\
07:58 & Two Y-up shifts in the fourth RK4 timeslot. & 7\\
08:00 & T, pyramid side view, circle, cone, sphere, apex, double dot, one-bit BST, $\delta\delta\phi$, Fibonacci phyllotaxis, blend and $\Psi$. & 8--9\\
08:04 & Crystal bifurcation of RGBA; A means inverse T. & 9--10\\
08:14 & Dichromatic operators and the remaining texture slots. & 11--12\\
08:16 & The common self-referential Klein-bottle SDF; optional downstream Bayer. & 12--14\\
08:23 & Operators themselves move and change in the LUT/BST at $\Psi$. & 14--16\\
08:30 & A GPU substrate, on-chip LUT and resident packed VRAM textures. & 18--19\\
08:33 & Hypothetical packing and operation-level complexity. & 20--21\\
08:44 & Zero-pointer implicit arrays and the Universal Spatial State Automaton designation. & 22--23\\
\bottomrule
\end{tabularx}}

\subsection{Symbols that remain distinct}
$\phi$ denotes the source's phase/angular slot. $\Phi$ names its log-encoded polar-radius operator. $\delta\delta\phi$ retains the source's double phase-differential name. $\Psi$ is the interval at which the field and its operator relations advance. The colon $:$ is the pinion double-dot coupling operator.

T has a geometric role in the primitive LUT and an inverse-related role in the RGBA construction. The notation uses $T_{\rm shape}$ and $T_{\rm transform}$ when those two positions need separate names. This keeps the source's uses visible without collapsing them into an ordinary opacity variable.
\clearpage

\section{Double pinion, mitosis and the one-bit split}
The double vertical slit is the twin-path element of the substrate. Each slit belongs to a pinion; the pinions drive the paired wavefront through the common surface. The Hadamard mitosis hinge names the splitting relationship. The later self-referential definition also gives the double pinion the role of keeping the operators on the Klein-bottle SDF as the operators move and change. \src{3--5, 14--16}

Let $W$ be an incoming wavefront and let
\begin{equation}
\bigl(W^{(0)},W^{(1)}\bigr)=\mathcal{H}_{\phi}[W]
\end{equation}
name its two hinge streams. The paired pinion action is written
\begin{equation}
W'=\mathcal{P}_{0,1}\bigl[W^{(0)},W^{(1)};\Klein,\LUT\bigr].
\end{equation}
$\mathcal{H}_{\phi}$ and $\mathcal{P}_{0,1}$ remain entries in the operator field; their phase, position and body may be changed by the same circulation in which they are used.

\subsection{The 07 / 54 instance}
The source uses the time components 07 and 54 as the two values in its one-bit illustration. Its parity table records the number of set bits. \src{4--5}

\begin{center}
\begin{tabular}{@{}lrrrr@{}}
\toprule
\textbf{Component} & \textbf{Decimal} & \textbf{Binary} & \textbf{Set bits} & \textbf{Parity}\\
\midrule
07 & 7 & \code{0111} & 3 & 1\\
54 & 54 & \code{110110} & 4 & 0\\
\bottomrule
\end{tabular}
\end{center}

For editing and integer evaluation, the workbench gives the two source expressions their own names:
\begin{align}
\operatorname{parity}(v)&=\operatorname{popcount}(v)\bmod 2,\\
\operatorname{even\_odd}(v)&=v\mathbin{\&}1.
\end{align}
Both give the displayed $1/0$ pair for the source's $7/54$ example. The parity value is carried into the branch and return expressions rather than treated as a detached measurement report.

\subsection{The split stays in the loop}
The two paths remain coupled by the pinion and the shared field. Their role is not exhausted by an initial split: the source subsequently reuses the paired structure as the central drive for returning state and changing operator definitions. The dichromatic channels carry these two paths into the texture state discussed in Section~\ref{sec:rgba}.
\clearpage

\section{Log-polar encoding and the FT wavefront}
The source introduces log-encoded polar vectors together with lowercase phi, and later names the log-encoded polar radius PHI explicitly. In the unified notation, a local wavefront address is written
\begin{equation}
q=\Phi[r,\phi;\LUT].
\end{equation}
This expression names the source's log-polar encoding operation. The LUT holds the relationship used to express the radius, phase and operator position. The encoding remains connected to the mutable LUT, rather than becoming a fixed Cartesian field laid beneath it. \src{3--4, 11--12}

The source associates logarithmic encoding with representing scaling and rotation as shifts in the encoded space. Within the workbench, these remain log-polar operator relationships. The log base, binning and numeric phase representation belong to the editable binding of $\Phi$.

\subsection{The source wavefront}
At the source label 07:56, the value 756 is assigned the wavefront FT vector
\begin{equation}
W_{756}=[0,2,0,1].
\end{equation}
The second component and the final one-bit component are then discussed in relation to the double slit and jitter. The formalization retains the vector as the stated source instance; the source does not supply a general timestamp-to-vector formula. \src{6}

Accordingly, the accompanying model contains
\begin{lstlisting}
"seed": {
  "timestamp_label": "07:56 / 756",
  "wavefront": [0, 2, 0, 1],
  "split_example": {"left": 7, "right": 54}
}
\end{lstlisting}

The wavefront enters the same recurrence as every later state. Its encoded form is passed into the split and hinge, then through the selected operator, kinematic slots, geometric divergence and RGBA state. Once a cycle has returned, the full preceding state becomes the next input.

\subsection{Encoding operators as well as values}
A log-polar entry can represent an operator position or relation as well as a wavefront value. This is the connection made explicit in the source's later formulation: jitter acts on the operators in the LUT around the Klein surface and changes their positions in the BST relation. Log-polar encoding therefore participates in operator motion, selection and recurrence, not only in a final display coordinate. \src{14--16}
\clearpage

\section{PSI, one-bit jitter and the fourth RK4 slot}
$\Psi$ is the interval of circulation and operator change. One-bit jitter is represented by $j_n\in\{0,1\}$, but its object is the entire operator-bearing state. The source's later construction assigns jitter an active role: it moves and changes the operators that will perform the next action. \src{8--9, 14--16}

Write the jitter-dependent operator update as
\begin{equation}
\widetilde{\LUT}_n=
\mathcal{J}_{\Psi}\bigl[\LUT_n,\State_n;j_n\bigr].
\end{equation}
The tilde names the operator field used during the current written cycle. The previous output is already contained in $\State_n$, so the change is self-referential. The $\Psi$ cycle and this choice of written ordering are exposed in \code{model/cycle.json}.

\subsection{Four timeslots, with the Y-up double event}
The source identifies four RK4 timeslots and places two Y-up shifts in the fourth slot. This relationship is written without separating the Y-up event from its kinematic action:
\begin{align}
k_1&=\mathcal{R}_1[W_n;\widetilde{\LUT}_n,\Psi,j_n],\\
k_2&=\mathcal{R}_2[W_n;\widetilde{\LUT}_n,\Psi,j_n],\\
k_3&=\mathcal{R}_3[W_n;\widetilde{\LUT}_n,\Psi,j_n],\\
\widehat{k}_4&=\mathcal{Y}_{\uparrow}
 \bigl[\mathcal{Y}_{\uparrow}[\mathcal{R}_4[W_n;\widetilde{\LUT}_n,\Psi,j_n]]\bigr],\\
W^{\rm kin}_n&=\mathcal{R}_{\rm RK4}[k_1,k_2,k_3,\widehat{k}_4;\Psi].
\end{align}
The four subscripts identify the source's timeslots. The Y-up displacement and vector derivative attach to their editable operator bindings. \src{7}

\subsection{Jitter is carried forward}
The next geometric and RGBA expressions contain the kinematic result. B carries jitter and RK4 history. The returned state is then available when the following interval changes the LUT again. The one-bit action consequently appears in the wavefront, operator body, operator position, routing context and history of the same cycle. \src{11--16}

The example session supplies an explicit, editable bit sequence for constructing traces. Jitter generation and the symbolic interval $\Psi$ retain their own named bindings.
\clearpage

\section{Geometric LUT, double dot and phyllotaxis}
At the source's 08:00 stage, the LUT contains T, a side view of the pyramid, a circle, a cone, a sphere and an apex. The apex is associated with the pinion double-dot product operator. The same stage introduces the one-bit BST, lowercase delta delta phi, Fibonacci phyllotaxis and the math blend operator at interval $\Psi$. \src{8--9}

\begin{equation}
\mathcal{G}=\{
 T_{\rm shape},\ \mathrm{pyramid}_{\rm side},\ \mathrm{circle},\
 \mathrm{cone},\ \mathrm{sphere},\ \mathrm{apex}\}.
\end{equation}
Every member is an SDF operator entry in the shared LUT. Its geometry is part of its action and its placement, rather than a separate list of visible objects.

\subsection{The pinion colon is its own operator}
For operands $U$ and $V$, write
\begin{equation}
U:_{\rm pinion}V
=\mathcal{C}_{:}[U,V;\Klein,\LUT].
\end{equation}
The colon retains the source's pinion-coupling name and its geometric relationship to the apex. Its operand types and coupling equation are editable in the operator body. The product is defined by that body and its position in the pinion relation.

\subsection{Divergence and blend are one process}
The phase-differential and phyllotaxis relationship is written
\begin{align}
D_n&=\mathcal{D}_{\delta\delta\phi}
 [\Phi_n,j_n;\widetilde{\LUT}_n,\Psi],\\
F_n&=\mathcal{F}_{\rm phyllotaxis}
 [D_n,\mathcal{G}_n;\mathcal{B}_n,\Psi],\\
U_n&=\mathcal{M}_{\rm blend}
 [W^{\rm kin}_n:_{\rm pinion}F_n,\ F_n,\ \mathcal{G}_n].
\end{align}
Here $F_n$ is both a divergence relation and an arrangement of operator positions. The source relates the arrangement to the Fibonacci/golden-angle spiral and later uses it when discussing neighboring execution work on the GPU. Its stated approximate spacing is $137.5^{\circ}$. \src{8--9, 19--21}

The six primitive entries, phase differential, coupling and blend are invoked within the same workbench cycle. Their changed bodies and locations are read from the current LUT. Editing one of these entries therefore alters an expression that participates in the returned state.
\clearpage

\section{The Klein bottle as the common operator field}
The source places the preceding operators inside a self-referential Klein-bottle SDF. The double pinion holds the operator flow to the surface. The return through the non-orientable relation connects output to input and carries a reversal of the route/parity relationship. \src{12--16}

Write the surface relation as
\begin{equation}
\Klein_n=\operatorname{SDF}_{\rm Klein}[\State_n].
\end{equation}
The named field is retained directly. Its numerical surface equation, coordinate chart and distance evaluation are exposed as the $\operatorname{SDF}_{\rm Klein}$ binding. The surface term carries the source's topology and operational roles into the recurrent expression.

\subsection{An operator is a situated definition}
A working operator record has the form
\begin{equation}
\Op_{i,n}=(\,\mathrm{body}_{i,n},\ f_{i,n},\ a_{i,n}\,),
\qquad \LUT_n[i]=\Op_{i,n},
\end{equation}
where $f$ is its SDF expression and $a$ its surface-position expression. The record's action is written
\begin{equation}
W'=\Apply[\Op_{i,n},W;\State_n].
\end{equation}
The field, location and body are all accessible to the next operator-change expression. This is the source's code/data unity written as a directly editable record structure.

\subsection{The pinion carries, the jitter changes}
In the written $\Psi$ cycle, each record is changed by
\begin{align}
\widetilde{\mathrm{body}}_{i,n}
 &=\mathcal{J}_{\rm body}
 [\mathrm{body}_{i,n},f_{i,n},\LUT_n,\State_n;j_n,\Psi],\\
\widetilde{a}_{i,n}
 &=\Pin_{\Klein}
 [a_{i,n},\State_n;j_n,\Psi],\\
\widetilde{f}_{i,n}
 &=\mathcal{J}_{\rm SDF}
 [f_{i,n},\widetilde{\mathrm{body}}_{i,n},\widetilde{a}_{i,n};\Klein].
\end{align}
These are named source relationships expressed as a working record update. The workbench keeps the body, field and surface-position dependencies intact in each evolving record.

The operator that changes a body is itself an entry in the same field. Its record is read from the preceding LUT when the cycle changes the LUT, preserving the self-referential dependency without copying an expanding expression into every operand.
\clearpage

\section{One-bit BST and zero-pointer routing}
The BST is the source's structural relationship between the SDF operators. The initial account describes binary paths through the field; its later optimization uses a zero-indexed implicit array whose child addresses follow directly from the parent index. \src{8--9, 13--15, 22}

\begin{equation}
\boxed{\operatorname{child}(i,b)=2i+1+b,\qquad b\in\{0,1\}.}
\end{equation}
Thus the left child is $2i+1$ and the right child is $2i+2$. An entry stores its body, field and position; separate left/right address fields are not part of this array relation.

For a route $b_0,b_1,\ldots,b_{d-1}$ starting at the root,
\begin{equation}
i_0=0,\qquad i_{k+1}=2i_k+1+b_k.
\end{equation}
The integer address selects a LUT entry whose current SDF record performs the action.

\subsection{Non-orientable reversal belongs to the route}
The source connects passage through the Klein-bottle neck with parity reversal and inversion of the execution path. A one-bit notation for that route relationship is
\begin{equation}
b'_k=b_k\oplus\sigma_n,
\end{equation}
where $\sigma_n$ names the neck/reversal event. The event's geometric evaluation remains connected to the Klein field. The example session supplies explicit reversal bits so that the written routing relation can be exercised and inspected. \src{15}

For example, the workbench route $[0,1]$ follows $0\to1\to4$. With reversal bit one, its directions become $[1,0]$ and it follows $0\to2\to5$. These are integer evaluations of the written array relation, not a new numerical formula for the Klein surface.

\subsection{The initial catalogue is editable}
The accompanying model stores the source terms in a 31-entry working catalogue. The array order is an editing choice, visible in \code{model/substrate.json}; the source does not assign this particular index sequence. Operator bodies can be replaced and further entries can be added. No fixed maximum depth or catalogue size is imposed by the workbench.

A referenced address without a supplied record stays a named \code{LUT\_read} expression. It is not redirected by modulo arithmetic to an unrelated occupied entry. This keeps the intended route visible while the definition is being developed.
\clearpage

\section{Dichromatic RGBA and A as inverse T}\label{sec:rgba}
The source's RGBA state develops from crystal bifurcation and then becomes dichromatic: two channels carry the paired pinion streams while the remaining texture roles carry the history and inverse relation. \src{9--12}

\begin{equation}
C_n=\bigl(R_n,G_n,B_n,A_n\bigr),
\qquad \boxed{A_n=T_n^{-1}.}
\end{equation}

\begin{tabularx}{\linewidth}{@{}P{17mm}Y@{}}
\toprule
\textbf{Channel} & \textbf{Role in the source definition}\\
\midrule
R & Log-polar state of one slit/pinion stream.\\
G & Log-polar state of the paired slit/pinion stream.\\
B & Phase/jitter history and the kinematic RK4 context.\\
A & Inverse T, carried with the dichromatic state.\\
\bottomrule
\end{tabularx}

\subsection{The inverse remains an inverse relation}
The source discusses inverse time, inverse transformation and inverse temperature as possible readings of inverse T. Its later texture description uses the inverse-transformation loop. The working model retains the chosen role as \code{T\_inverse\_role}, initially \code{transformation}, and keeps the inverse as a named action on the T expression. \src{10--12}

The expression is not replaced by a clipped display-opacity number. In the current workbench it is built as
\begin{equation}
T_n=\mathcal{T}_{\rm transform}[U_n,\operatorname{Crystal}_n],
\qquad A_n=\mathcal{I}_T[T_n].
\end{equation}
Changing the body of \code{inverse\_T} or the T transformation changes this same state path. The numerical meaning of the inverse is attached to that chosen T binding.

\subsection{Crystal, history and return}
Crystal bifurcation supplies the RGBA expression from the geometric/kinematic result. The two active channels remain paired. B carries the previous history into the new history expression alongside the one-bit event and four kinematic slots. A retains the inverse-T expression. All four then return through the pinion and Klein surface.

The source's mentions of inverse-related stabilization remain part of its A-channel description. In this edition that behavior lives in the editable inverse-T operator, rather than in a separate workbench filter that discards or clamps the evolving state.
\clearpage

\section{The complete PSI circulation}
The following is one written ordering of the source's interacting relationships. It is the order exposed in the accompanying cycle file, so the document and executable expression construction can be edited together.

{\small
\begin{tabularx}{\linewidth}{@{}P{10mm}P{39mm}Y@{}}
\toprule
\textbf{Step} & \textbf{Action} & \textbf{State carried onward}\\
\midrule
1 & Change the operator field & Previous whole state, LUT and jitter change operator bodies, SDFs and pinion-held surface positions.\\
2 & Encode through PHI & The current input is expressed with the current log-polar LUT.\\
3 & Split and hinge & The double-slit bit relation and Hadamard mitosis produce the paired stream with parity context.\\
4 & Select through the BST & The current route and reversal relation select a situated operator record.\\
5 & Advance four RK4 slots & The selected result enters all four slots; Y-up acts twice in the fourth.\\
6 & Couple geometry and divergence & Primitive SDFs, $\delta\delta\phi$, phyllotaxis, colon coupling and blend act in one expression.\\
7 & Form the RGBA crystal & R/G carry the pair; B carries history; A carries inverse T.\\
8 & Return through the surface & The double pinion and non-orientable return form the next whole state and its next input.\\
\bottomrule
\end{tabularx}}

\begin{align}
\widetilde{\LUT}_n&=\mathcal{J}_{\Psi}[\LUT_n,\State_n;j_n],\\
W'_n&=\mathcal{E}_{\Psi}[\State_n;\widetilde{\LUT}_n,j_n],\\
C'_n&=\mathcal{C}_{\rm RGBA}[W'_n,\State_n;\widetilde{\LUT}_n],\\
\State_{n+1}&=\mathcal{R}_{\Klein}
 [\State_n,W'_n,C'_n;\widetilde{\LUT}_n],\\
\operatorname{Input}_{n+1}&=\State_{n+1}.
\end{align}
The symbols $\mathcal{E}$, $\mathcal{C}$ and $\mathcal{R}$ collect the named actions in the table; they do not introduce a second numerical engine. The full returned state includes the changed operator field, so the next input and the next acting definitions remain connected. \src{13--17}

An optional readout is written $Y_n=\mathcal{B}_{\rm Bayer}[\State_n]$. It is evaluated after state return. Its expression does not become a required input to the next cycle.
\clearpage

\section{Self-reference, operator change and expression}
The source's shift from graphics to computing is the shift from transporting only values to transporting and changing the operators that act on them. Its double pinion, surface and jitter are therefore part of the computing definition itself. \src{14--16}

A cycle may change an operator body, its surface position and the field expression through which it is selected or applied. In the workbench, all three are present in the operator record. Every application refers to the current record. The old record remains in the expression history because the new record may depend on it.

\subsection{State becomes instruction}
The recurrent relationship can be read in either direction:
\begin{equation}
\text{operator field}\ \longrightarrow\ \text{returned state}
\ \longrightarrow\ \text{changed operator field}.
\end{equation}
The source describes this as data becoming instruction and as a self-contained evolutionary processing layer. The change is not confined to a side channel: the operator's changed definition is the definition referenced by the next application. \src{15--16}

\subsection{The expression vocabulary of the whole}
The source's Universal Spatial State Automaton designation places several workloads in the same language. \src{22--23}

\begin{tabularx}{\linewidth}{@{}P{37mm}Y@{}}
\toprule
\textbf{Expression} & \textbf{Its source representation}\\
\midrule
Conditional logic & A spatial branch choosing a path through the operator structure.\\
Repeated computation & The returned state circulating through the Klein relation.\\
Changing a program & Jitter and feedback changing the operator-bearing LUT.\\
Neural weights and activations & Interference/resonance values carried by the dichromatic LUT.\\
Database state and keys & Encoded log-polar coordinate/state relationships.\\
Downstream observation & The optional Bayer or other output expression attached to returned state.\\
\bottomrule
\end{tabularx}

These are the source's expression mappings, retained together as the system's design vocabulary. The package keeps these terms in the same situated-operator description. Further operators can be registered using the same body, field and surface-position record, and can then be routed to through the same implicit address relation.
\clearpage

\section{GPU substrate and resident texture organization}
The source identifies the owner's target as an \textbf{RTX 5070 Ti laptop edition with 12 GB VRAM}. Its GPU discussion places a compact one-bit LUT and changing operator state in the on-chip layer, with the packed texture state in VRAM. These are the source's target and residency objectives recorded in \code{model/gpu\_target.json}. \src{18--19}

\begin{tabularx}{\linewidth}{@{}P{32mm}Y@{}}
\toprule
\textbf{Location} & \textbf{Role in the source architecture}\\
\midrule
On-chip / cache layer & The compact one-bit LUT, binary routing and frequently changing operator context.\\
VRAM texture pool & The dense encoded state, texture records and circulating history.\\
Host exchange & Initial loading, explicit updates and requested downstream exchange.\\
\bottomrule
\end{tabularx}

\subsection{Residency belongs to the execution plan}
The source's intended operation is to load the packed working state and reuse it through self-reference, without making a hot-swap of textures a required step of each circulation. The two pinion channels keep their dichromatic roles; B and A retain the history and inverse-T roles in the overall texture definition. BC5 is the packing name used in the source's R/G discussion. \src{19--22}

The workbench represents the GPU layout as an editable target record. Device allocation, on-chip placement and texture bindings are editable entries for the device implementation, alongside the recurrent-state and operator records.

\subsection{Work follows the operator arrangement}
The source uses Fibonacci phyllotaxis to organize neighboring work and discusses neighboring surface regions being processed together. In the unified design, a group follows the same $\Psi$ cycle while referring to the operator field and encoded state appropriate to that group. \src{19--21}

\begin{lstlisting}
load encoded substrate and operator LUT
repeat a chosen number of PSI intervals:
    change situated operator definitions
    apply pinion / BST / kinematic / geometry cycle
    return the RGBA state through the Klein relation
read downstream output when requested
\end{lstlisting}

This is the device-side expression plan carried by the source architecture. The accompanying Python workbench exports the corresponding expression graph, live operator field and target records for continued implementation work.
\clearpage

\section{Packing scenarios and operation-level cost}
The source explicitly introduces its packing discussion as hypothetical. Its table uses a 12 GB pool label and reports the following approximate counts. They are recorded here as the source's packing scenarios. \src{20--22}

{\small
\begin{tabularx}{\linewidth}{@{}YP{21mm}P{36mm}@{}}
\toprule
\textbf{Source representation} & \textbf{Bytes / element} & \textbf{Source approximate count}\\
\midrule
RGBA32F & 16 & 805,000,000\\
RGBA8 & 4 & 3,220,000,000\\
Dichromatic BC5 & 1 & 12,880,000,000\\
Implicit-tree operational-state scenario & --- & 24--48 billion\\
\bottomrule
\end{tabularx}}

\subsection{A directly editable packing calculation}
For a selected byte budget $M$ and stated element width $b$, the calculator evaluates
\begin{equation}
N(M,b)=\left\lfloor\frac{M}{b}\right\rfloor.
\end{equation}
Its default input is $M=12\times2^{30}=12{,}884{,}901{,}888$ bytes, which reproduces the scale of the source's rounded figures. The byte budget is an input rather than a discovered property of a device.

The source's later 24--48 billion range is stored as a separate named scenario. It does not supply a bytes-per-operational-state definition for that range, so that field remains open in the packing record instead of being backfilled with an invented record size.

\subsection{Costs refer to their named operation}
The source assigns $O(1)$ to the static log-polar lookup and to changing the shared one-bit operator control. It assigns $O(\log_2N)$ to the implicit binary route. These labels are attached to those specific operations in \code{model/packing\_scenarios.json}. \src{21--22}

A route supplied as $d$ bits takes $d$ applications of the child relation in the workbench. Its symbolic cycle also constructs records for its operator entries and for the complete wavefront/geometry/return expression. Those counts are emitted as \code{operator\_count} and \code{expression\_node\_count}, so the actual workbench run has its own inspectable record rather than borrowing the source's hardware scenarios.

\begin{lstlisting}
python run.py packing
python run.py packing --bytes 12000000000
\end{lstlisting}
\clearpage

\section{Editable bindings and the living definition}
The source defines a connected architecture, but leaves several numerical functions unnamed beyond their operator terms. Those terms remain first-class entries in the workbench. The following binding positions collect them for editing rather than substituting unrelated numerical behavior.

{\small
\begin{tabularx}{\linewidth}{@{}P{35mm}YP{15mm}@{}}
\toprule
\textbf{Binding} & \textbf{Value or law to enter} & \textbf{Source}\\
\midrule
Klein SDF & Surface equation, coordinate convention and field evaluation. & 12--16\\
PHI encoding & Log base, phase coding and numeric precision. & 3--4, 11\\
Hadamard / pinion & The phase-dependent hinge and pinion action. & 3--5\\
Double dot $:$ & Operand types and the coupling equation. & 8\\
$\delta\delta\phi$ / blend & The phase-differential law and blend parameters. & 8--9\\
RK4 / Y-up & The vector derivative and the value of a Y-up action. & 7\\
$\Psi$ / jitter & The numeric interval and bit-generation relation. & 8--9, 14--16\\
Inverse T & The inverse for the selected T role. & 9--12\\
Device packing & Record layout, allocation and device bindings. & 18--22\\
\bottomrule
\end{tabularx}}

\subsection{Change the definition where it is used}
The operator catalogue lives in \code{model/substrate.json}. Each entry has a name, implicit index, body and field description. Bodies are JSON expressions. The runtime provides \code{set\_body} for replacing one, and \code{add\_operator} for registering another situated definition. The next circulation uses the resulting record in the same source-native action path.

\begin{lstlisting}
# After loading a Kernel instance named kernel:
kernel.set_body("double_dot", {
    "symbol": "My_double_dot_body",
    "operands": ["left_pinion", "right_pinion", "local_SDF"]
})
kernel.add_operator("my_operator", "My SDF operator", {
    "symbol": "My_operator_definition"
})
\end{lstlisting}

These expressions are authoring material. The included workbench executes symbolic expression construction and operator-record evolution; a numerical implementation supplies the numerical meaning of the named body. The full returned state and operator dependencies remain available in the exported graph for that work.
\clearpage

\section{Working with the accompanying files}
The ZIP contains one attributed edition and its editable materials. The PDF, manuscript, model and code all use this unified definition. The original 24-page discussion remains unchanged in the source folder.

\begin{tabularx}{\linewidth}{@{}P{37mm}Y@{}}
\toprule
\textbf{Location} & \textbf{Working material}\\
\midrule
\code{docs/} & This PDF, editable LaTeX and Markdown manuscripts, operator dictionary.\\
\code{model/} & Substrate catalogue, cycle, numerical binding positions, GPU target and source packing scenarios.\\
\code{src/dawnwood/} & Expression graph, concrete one-bit/array relations and unified symbolic circulation.\\
\code{examples/} & Source seed, editable session inputs and a body-edit example.\\
\code{tests/} & Tests for graph identity, routing, operator changes, feedback, history and optional readout.\\
\code{results/} & Generated run, test report and reproduction record.\\
\code{source/} & Original PDF and extracted page-by-page text.\\
\code{author.json} & Tom Klootwijk's supplied attribution and identifiers.\\
\bottomrule
\end{tabularx}

\subsection{Start a circulation}
From the extracted package directory, use Python 3.10 or later:
\begin{lstlisting}
python run.py catalogue
python run.py run --steps 16 --out work/session
python run.py inspect work/session/snapshot.json \
    --operator double_dot
\end{lstlisting}
The trace records the selected implicit route, one-bit inputs, state identities and expression counts. The snapshot contains the live operator records and complete shared expression graph. No external Python packages are required for these commands.

\subsection{Continue, edit and compare}
\begin{lstlisting}
python run.py run --resume work/session/snapshot.json \
    --steps 16 --out work/continued
python run.py run --body-edit examples/body_edit.json \
    --steps 16 --out work/edited
python tools/run_tests.py
python tools/build_pdf.py
python tools/verify_manifest.py
\end{lstlisting}
The step count is a chosen work session length. The workbench applies no fixed upper step count, chart-size cap, fixed opcode whitelist, opacity clamp, radius gate or automatic route aliasing. Bit and index checks express their declared types. The editable cycle and operator records carry the behavior.
\clearpage

\section*{Source and edition record}
\addcontentsline{toc}{section}{Source and edition record}
\textbf{Source}\par
\code{double-slit-theory.pdf}, 24-page exported discussion supplied with the request. The original is included at \code{source/double-slit-theory.pdf}. The page references in this document point only to that source.

\textbf{Working formalization}\par
Version 0.2 organizes the discussion's operator relationships into a single recurrence and an editable expression model. Mathematical symbols, operator-record layout, initial catalogue order, example bit sequences and the written cycle order are authoring notation introduced by this edition. They are exposed directly in the accompanying files.

\textbf{Source navigation}\par
\begin{tabularx}{\linewidth}{@{}P{25mm}Y@{}}
\toprule
\textbf{Pages} & \textbf{Material}\\
\midrule
3--7 & Double pinion, log-polar phi, parity, FT wavefront, jitter and fourth RK4 timeslot.\\
8--12 & Primitive LUT, double dot, phase differential, Fibonacci phyllotaxis, RGBA and inverse T.\\
12--17 & Common Klein-bottle SDF, operator BST, optional Bayer and self-referential operator change.\\
18--21 & GPU target, on-chip/VRAM arrangement and hypothetical packing/cost scenarios.\\
22--23 & Implicit zero-pointer array and the Universal Spatial State Automaton expression language.\\
\bottomrule
\end{tabularx}

\vspace{8mm}
\begin{dwbox}[Dawnwood Interactive]
\textbf{Tom Klootwijk}\par
Identifier: \textbf{NL200678942}\par
Date of birth: \textbf{10-07-1990}\par
Telephone: \textbf{+31655954068}\par
\vspace{3mm}
\textbf{Unified Substrate Definition}\par
Version 0.2 \enspace | \enspace 16 September 2026
\end{dwbox}
\vfill
{\sffamily\color{DWGreen}\large The definition returns to itself.}
\end{document}
